% !TeX root = ../drafts/c-flock-protocol.tex
\section{The Flock protocol}\label{annex:flock}

Flock~\cite{Flock26} proves many executions of one Boolean circuit at once, in our case the BLAKE2s compression.

\subsection{Statement and commitment}\label{flock:statement}

One BLAKE2s compression has a Boolean R1CS witness block of $2^{\kblock}$ bits, where $\kblock=14$. Its fixed matrices are $A_0,B_0,C_0\in\Ftwo^{2^{\kblock}\times 2^{\kblock}}$, with $C_0=I$. The witness includes the circuit inputs, outputs, AND-gate values, addition carries and one constant-one position; linear wires are substituted away. Within-block coordinates are the low-order coordinates and batch coordinates are the high-order coordinates throughout this annex.

For $2^{\kbatch}$ independent compressions, the witness is:
\[
z:\cube{\kbool}\longrightarrow\Ftwo,\qquad \kbool=\kblock+\kbatch,
\]
and use
\[
A=I_{2^{\kbatch}}\otimes A_0,\qquad B=I_{2^{\kbatch}}\otimes B_0,\qquad C=I_{2^{\kbool}}
\]
Writing $a=Az$, $b=Bz$ and $c=Cz=z$, the R1CS statement is:
\begin{equation}\label{flock:eq:r1cs}
a(u)b(u)+c(u)=0\qquad\text{for every }u\in\cube{\kbool}
\end{equation}

We also enforce the position $512$ of every block to be $1$ (preventing the all-zero witness):
\begin{equation}\label{flock:eq:const-one}
z(0,0,0,0,0,0,0,0,0,1,0,0,0,0,\,t)=1\qquad\text{for every }t\in\cube{\kbatch}
\end{equation}

This second condition is enforced by lincheck (\S\ref{flock:lincheck}).

Set $\kskip=6$ and $\nflock=\kbool-\kskip$. Identify $i\in\cube{\kskip}$ with the integer whose binary expansion has the same low-order coordinates. The prover packs each group of $2^{\kskip}=64$ witness bits into one $\K$-element,
\[
\qflock:\cube{\nflock}\longrightarrow\K,\qquad \qflock(v)=\sum_{i=0}^{63}z(i,v)x^i,
\]
and commits to $\qflock$ (as a stacked multilinear of \S\ref{sec:stacking}). $\qflock$ will eventually receive opening claims both from Flock and from the rest of the VM interactions (see \S\ref{sec:tab-blake2s}).


\subsection{Zerocheck}\label{flock:zerocheck}

Flock zerochecks \eqref{flock:eq:r1cs} (Definition~\ref{def:zerocheck}), but replaces the first $\kskip$ rounds by one univariate round~\cite{Gru24}. Split $u=(i,v)$ with $i\in\cube{\kskip}$ and $v\in\cube{\nflock}$. Let $\fembed:\F_{2^8}\hookrightarrow\E$ be the subfield embedding, with a byte interpreted in the polynomial basis, and define
\[
\skipdom=\{\fembed(0),\ldots,\fembed(63)\},\qquad \skipcoset=\fembed(64)+\skipdom=\{\fembed(64),\ldots,\fembed(127)\}.
\]
Thus $\skipdom$ is a $64$-point additive subspace and $\skipcoset$ is its disjoint coset, in the additive-NTT ordering of Annex~\ref{annex:pcs}.

Write $\lag_i$ for the Lagrange basis of $\skipdom$: $\lag_i(\fembed(j))=[i=j]$. Extend $a$ by
\[
  \qext a(Y,v)\;=\;\sum_{i=0}^{63}\lag_i(Y)\,\mle a(i,v),\qquad \qext a:\E\times\E^{\nflock}\to\E,
\]
the multilinear extension in $v$, the $64$-point interpolation in $Y$. So $\qext a$ has degree $63$ in $Y$, not one, and it agrees with $a$ on $\skipdom\times\cube{\nflock}$. Define $\qext b$ and $\qext c$ likewise. Arguments come in that order everywhere below: the skip argument first, then the remaining within-block coordinates, then the batch ones. 

The zerocheck point $r \in \E^{\nflock}$ is defined as follows:
\begin{itemize}
  \item Its first seven coordinates $r_0, \dots, r_6$ are fixed to $\fembed(\mathtt{0xf7}),\fembed(\mathtt{0x53}),\fembed(\mathtt{0xb5})$ and $g_0^{2^j}/(1+g_0^{2^j})$ for $0\leq j<4$, where $g_0\in\E$ is public. Lemma~\ref{lem:fixed-zerocheck} applies.
  \item The remaining coordinates $r_7,\dots,r_{\nflock-1}$ are sampled at random by the verifier.
\end{itemize}

Define
\begin{align}
P^{AB}(Y)&=\sum_{v\in\cube{\nflock}}\eq(r,v)\,\qext a(Y,v)\qext b(Y,v),\label{flock:eq:pab}\\
P^C(Y)&=\sum_{v\in\cube{\nflock}}\eq(r,v)\,\qext c(Y,v).\label{flock:eq:pc}
\end{align}
Both lie in $\E[Y]$: $P^{AB}$ has degree at most $126$, $P^C$ degree below $64$. If \eqref{flock:eq:r1cs} holds, then $P^{AB}(h)+P^C(h)=0$ for every $h\in\skipdom$. The prover sends $P^{AB}|_{\skipcoset}$ and $P^C|_{\skipcoset}$, $64$ values each. The latter determines $P^C$; the assumed zero values on $\skipdom$ and the transmitted values of $P^{AB}+P^C$ on $\skipcoset$ determine their degree-below-$128$ sum $P^{AB}+P^C$. Finally, the verifier deduces $P^{AB} = (P^{AB}+P^C) - P^C$.

The verifier then samples $\zskip\in\E$ and reconstructs the two claimed values
\[
v_C \qeq P^C(\zskip),\qquad v_{AB} \qeq P^{AB}(\zskip).
\]

By Fact~\ref{fact:eq} and $c=z$, the first value corresponds to the claim:
\begin{equation}\label{flock:eq:c-claim}
\qext z(\zskip,r) \qeq v_C.
\end{equation}
Thus $C$ takes no further sumcheck rounds and needs no lincheck.

For $P^{AB}$, the parties run sumcheck over the remaining $\nflock$ Boolean variables in \eqref{flock:eq:pab}. At the end the prover sends $v_a,v_b\in\E$, which lets the verifier conclude the sumcheck, with the new claims:
\[
\qext a(\zskip,\rho)=v_a,\qquad \qext b(\zskip,\rho)=v_b,
\]
where $\rho\in\E^{\nflock}$ is the vector of sumcheck challenges.

\subsection{Lincheck and the circuit walk}\label{flock:lincheck}

Lincheck reduces the two zerocheck claims to claims on $z$. The block structure is what keeps it cheap: the matrices repeat one BLAKE2s circuit, so lincheck runs over a single block, not over the batch. Split $\rho=(\rho_{\mathrm{in}},\rho_{\mathrm{out}})$, where $\rho_{\mathrm{in}}\in\E^{\kblock-\kskip}$ selects the remaining within-block coordinates and $\rho_{\mathrm{out}}\in\E^{\kbatch}$ selects a batch instance. Let $\qext A_0$ extend the row index of $A_0$, interpolated over $\skipdom$ in its low $\kskip$ coordinates and multilinear in the other $\kblock-\kskip$, the column index staying Boolean.

\begin{lemma}[Block diagonality folds the batch into $z$]\label{lem:flock-block-diagonal}
For $a=Az$ with $A=I_{2^{\kbatch}}\otimes A_0$, at every point $(\zskip,\rho_{\mathrm{in}},\rho_{\mathrm{out}})\in\E\times\E^{\kblock-\kskip}\times\E^{\kbatch}$,
\[
\qext a(\zskip,\rho_{\mathrm{in}},\rho_{\mathrm{out}})=\sum_{j\in\cube{\kblock}}\qext A_0(\zskip,\rho_{\mathrm{in}},j)\,\mle z(j,\rho_{\mathrm{out}}),
\]
and likewise for $b=Bz$ with $B_0$ in place of $A_0$.
\end{lemma}

\begin{proof}
Block diagonality gives
\[
a(k,t)=\sum_{j\in\cube{\kblock}}A_0(k,j)\,z(j,t)\qquad\text{for }k\in\cube{\kblock},\;t\in\cube{\kbatch},
\]
so, splitting the row as $k=(i,k')$ with $i\in\cube{\kskip}$ and $k'\in\cube{\kblock-\kskip}$ and expanding every extension by Fact~\ref{fact:eq},
\begin{align*}
\qext a(\zskip,\rho_{\mathrm{in}},\rho_{\mathrm{out}})&=\sum_{i\in\cube{\kskip}}\lag_i(\zskip)\,\mle a(i,\rho_{\mathrm{in}},\rho_{\mathrm{out}})\\
&=\sum_{i\in\cube{\kskip}}\lag_i(\zskip)\sum_{k'\in\cube{\kblock-\kskip}}\eq(\rho_{\mathrm{in}},k')\sum_{t\in\cube{\kbatch}}\eq(\rho_{\mathrm{out}},t)\,a\bigl((i,k'),t\bigr)\\
&=\sum_{i,k'}\lag_i(\zskip)\,\eq(\rho_{\mathrm{in}},k')\sum_{j\in\cube{\kblock}}A_0\bigl((i,k'),j\bigr)\underbrace{\sum_{t\in\cube{\kbatch}}\eq(\rho_{\mathrm{out}},t)\,z(j,t)}_{=\;\mle z(j,\rho_{\mathrm{out}})}\\
&=\sum_{j\in\cube{\kblock}}\Bigl(\underbrace{\sum_{i,k'}\lag_i(\zskip)\,\eq(\rho_{\mathrm{in}},k')\,A_0\bigl((i,k'),j\bigr)}_{=\;\qext A_0(\zskip,\rho_{\mathrm{in}},j)}\Bigr)\mle z(j,\rho_{\mathrm{out}}).
\end{align*}
\end{proof}

The verifier samples $\lcab\in\E$ and combines the two identities of Lemma~\ref{lem:flock-block-diagonal} with target $\lcab v_a+v_b$, adding the constant-position check \eqref{flock:eq:const-one} at weight $\lcab^2$. Lincheck proves
\begin{equation}\label{flock:eq:lincheck}
\lcab v_a+v_b+\lcab^2\qeq\sum_{j\in\cube{\kblock}}\Bigl(\lcab\qext A_0(\zskip,\rho_{\mathrm{in}},j)+\qext B_0(\zskip,\rho_{\mathrm{in}},j)+\lcab^2[j=512]\Bigr)\mle z(j,\rho_{\mathrm{out}}).
\end{equation}

The parties prove it by a partial sumcheck, only the $8$ rounds that bind the high column coordinates, leaving the low $\kskip = \kblock - 8 = 6$ unfolded. Write $\rho'_{\mathrm{in}}$ for those challenges, in coordinate order, which is the reverse of their round order, and $R_{\mathrm{lc}}\in\E$ for the claim they leave.

The prover then sends the $64$ unfolded slices, $\mle z(\cdot,\rho'_{\mathrm{in}},\rho_{\mathrm{out}})$:
\[
\zpartial=(s^{AB}_0,\ldots,s^{AB}_{63})\in\E^{64}.
\]

The verifier checks that sumcheck's terminal identity, which is \eqref{flock:eq:lincheck} with $\mle z(j,\rho_{\mathrm{out}})$ replaced by the folded column weights $w_{\mathrm{col}}(i,j')=s_i^{AB}\,\eq(\rho'_{\mathrm{in}},j')$:
\begin{equation}\label{flock:eq:lincheck-terminal}
R_{\mathrm{lc}}\qeq\sum_{j\in\cube{\kblock}}\Bigl(\lcab\qext A_0(\zskip,\rho_{\mathrm{in}},j)+\qext B_0(\zskip,\rho_{\mathrm{in}},j)+\lcab^2[j=512]\Bigr)w_{\mathrm{col}}(j).
\end{equation}
Both parties evaluate that sum by walking the BLAKE2s circuit backwards, one circuit against the roughly $89$ million nonzero entries the two matrices hold.

Once $\zpartial$ passes that check it is the $AB$ claim:
\begin{equation}\label{flock:eq:ab-claim}
\mle z(i,\rho'_{\mathrm{in}},\rho_{\mathrm{out}})=s_i^{AB},\qquad 0\leq i<64,
\end{equation}
and ring switching binds all $64$ of those values at once.

\paragraph{Ring switching.} Flock has now reduced every BLAKE2s constraint to two claims on $z$: the $C$ claim \eqref{flock:eq:c-claim} and the $AB$ claim \eqref{flock:eq:ab-claim}. Ring switching (Annex~\ref{annex:rs}) turns each into an MLE-friendly weighted claim on the packed $\qflock$ (Definition~\ref{def:ipcs}), under one shared $\Ftwo$-linear map from six challenges. Each is $64$ claims over the non-skipped coordinates: for $AB$ they are $\zpartial$, already sent and checked, so only the $C$ vector $s_0^C,\ldots,s_{63}^C\in\E$ is sent here, tied to $v_C$ by $v_C\qeq\sum_i\lag_i(\zskip)s_i^C$. The two weighted claims then join every other claim on the stacked witness, and one WHIR opening (Annex~\ref{annex:pcs}) discharges the combination.

\subsection{Protocol and soundness summary}\label{flock:summary}

\begin{protocol}[The Flock reduction, as implemented]\label{flock:protocol}
After the public statement, announced sizes, and stacked witness commitment have been bound:
\begin{enumerate}
\item The verifier constructs the zerocheck equality point from seven fixed friendly coordinates and sampled outer coordinates. The $\kskip$ variables replaced by the univariate skip consume no equality challenges.
\item The prover sends $P^{AB}|_{\skipcoset}$ and $P^C|_{\skipcoset}$, $64$ values each. The verifier samples $\zskip$, derives the $C$ claim, and reconstructs the initial $AB$ claim.
\item The parties run $\nflock$ quadratic zerocheck rounds. The prover sends $v_a,v_b$ before the verifier samples lincheck's batching challenge.
\item Lincheck samples $\lcab$, runs $\kblock-\kskip=8$ quadratic rounds, receives $\zpartial$, and checks \eqref{flock:eq:lincheck-terminal}, which leaves $\zpartial$ as the $AB$ claim.
\item Ring switching reuses $\zpartial$, receives and checks only the $C$ vector, samples its six shared map challenges, and turns the two Boolean claims into claims on $\qflock$.
\item The two claims join every other claim on the stacked witness, and the single PCS opening proves their random linear combination.
\end{enumerate}
\end{protocol}

For soundness, list binding first implies that, except with the PCS error, an accepting transcript has some member of $\polyset_{\mathsf{cm}}$ satisfying every opened claim (Definition~\ref{def:listbinding}). Fix such a member. Decomposing each of its $\K$-coordinates in the fixed $\Ftwo$-basis uniquely determines candidate bit multilinears and hence a candidate $z$ for that member, and $r$, $\zskip$ and every later challenge are sampled after the commitment that fixed it.

Suppose that candidate violates an R1CS row \eqref{flock:eq:r1cs}, at $u=(i,v)$ say. By Fact~\ref{fact:eq}, $P^{AB}(\fembed(i))+P^C(\fembed(i))$ is the extension at $r$ of the Boolean multilinear $v'\mapsto a(i,v')b(i,v')+c(i,v')$, which the violated row makes nonzero, and the seven fixed coordinates of $r$ carry $\Ftwo$-independent equality weights, so they cannot erase that violation and the sampled coordinates catch what survives: $P^{AB}+P^C$ is nonzero at $\fembed(i)$ except with probability $(\nflock-7)/|\E|$ (Lemma~\ref{lem:fixed-zerocheck}). Its $64$ zeros on $\skipdom$ are never tested, only assumed, so let $\hat P^C$ be the degree-below-$64$ polynomial the verifier interpolates from the received $P^C|_{\skipcoset}$ and $\hat P^{AB}$ the degree-below-$128$ one it then deduces, its degree being bounded only by the $64+64$ interpolation points, not by the honest $126$. If the received values are not those of the true $P^C$ then $\hat P^C\neq P^C$, so $v_C\neq P^C(\zskip)$ except with probability $63/|\E|$ (Corollary~\ref{cor:idtest}) and the $C$ claim \eqref{flock:eq:c-claim} is false about the candidate. Otherwise $\hat P^C=P^C$, and because $\hat P^{AB}+\hat P^C$ vanishes on all of $\skipdom$ by construction while $P^{AB}+P^C$ does not, $\hat P^{AB}\neq P^{AB}$ and the same corollary makes $v_{AB}$ false except with probability $127/|\E|$. A false $v_{AB}$ enters the $\nflock$ zerocheck rounds, which reject except with probability $2\nflock/|\E|$ (Fact~\ref{fact:sumcheck}, one degree-two cofactor per round), unless $v_a$ and $v_b$ are themselves false about the candidate. Lincheck binds $v_a,v_b$ to that same candidate and pins the constant-one position \eqref{flock:eq:const-one}, so either way an accepting transcript is left carrying a false claim on $z$.

Lemma~\ref{lem:sz} also does the list collapse: after the seven friendly coordinates are fixed, Lemma~\ref{lem:fixed-zerocheck} keeps the difference of two distinct Boolean candidates from vanishing identically, so the pair collides at the point of the $C$ claim with probability at most $(\kbool+50)/|\E|$, that point having degree $63$ in $\zskip$ and degree one in each of the $\kbool-13$ free coordinates of $r$.

Ring switching binds each resulting bit claim to $\qflock$ with error below $2^{-160}$ before the PCS error (Annex~\ref{annex:rs}). Because two claims share the map, a union bound gives total ring-switch error below $2^{-159}$. The verifier samples each challenge only after it has received the messages that challenge batches. The construction is a succinct argument, not a zero-knowledge proof.

The prover's dominant work is linear in the total number of Boolean witness bits. The univariate skip crosses from bits to $\E$ once, its $64$-point row extensions use byte lookup tables, and lincheck's fixed-matrix work depends on one BLAKE2s circuit after the batch coordinates are folded. Proof length and verifier work grow only logarithmically with the batch, apart from the fixed circuit walk and the final PCS queries.
